\lecture{19}{\texorpdfstring{PL-неравенство\\ и конечное время}{PL-неравенство и конечное время}}

\sourcevideo
    {https://www.youtube.com/playlist?list=PLOzRYVm0a65fhKDS8cYIC7YCuVGJ4FOJx}
    {Week 6: Lecture 19: Advanced Results on PL Inequality: Part 1}

Масштабирование положительным скаляром сохраняет геометрические
траектории непрерывной системы и изменяет только скорость движения по
ним. Для ньютоновского поля это позволяет получить фиксированное время
сходимости без равномерной константы сильной выпуклости. Во второй
части лекции доказывается геометрическое следствие неравенства
Поляка--Лоясиевича: квадратичный рост относительно всего множества
глобальных минимумов.

\section{Критерии времени и масштабирование поля}

Пусть \(V\) положительно определена относительно равновесия
\(x^\star\). Если
\[
    \dot V\le-cV^\alpha,
    \qquad
    c>0,
    \qquad
    0<\alpha<1,
\]
то равновесие достигается за конечное время и
\[
    T(x_0)
    \le
    \frac{V(x_0)^{1-\alpha}}{c(1-\alpha)}.
\]
При экспоненциальной оценке \(\dot V\le-cV\) функция обычно лишь
стремится к нулю. Степень \(\alpha<1\) не позволяет убыванию стать
слишком медленным около равновесия.

Если
\[
    \dot V
    \le
    -c_1V^{\alpha_1}
    -c_2V^{\alpha_2},
\]
где
\[
    c_1,c_2>0,
    \qquad
    0<\alpha_1<1<\alpha_2,
\]
то время установления равномерно ограничено:
\[
    T(x_0)
    \le
    \frac1{c_1(1-\alpha_1)}
    +
    \frac1{c_2(\alpha_2-1)}.
\]
Степень выше единицы отвечает за быстрое движение вдали от решения,
а степень ниже единицы --- за точное достижение после входа в его
окрестность.

Пусть система
\[
    \dot x=F(x)
\]
имеет устойчивое равновесие и \(a(x)>0\) вне него. Система
\[
    \dot x=a(x)F(x)
\]
задаёт те же непараметризованные фазовые траектории и отличается
перепараметризацией времени. Чтобы получить конечно- или
фиксированно-временной результат, необходимо отдельно проверить
определение поля в равновесии, существование решений, устойчивость и
нужное неравенство для функции Ляпунова.

После явной дискретизации это геометрическое совпадение теряется.
Например, фиксированно-временной градиентный поток формально приводит
к рекурсии
\[
\begin{aligned}
    x^{k+1}
    ={}&
    x^k
    -\eta
    \frac{\nabla f(x^k)}{
        \lVert\nabla f(x^k)\rVert^{\frac{p-2}{p-1}}
    }
    \\
    &-\eta
    \frac{\nabla f(x^k)}{
        \lVert\nabla f(x^k)\rVert^{\frac{q-2}{q-1}}
    },
\end{aligned}
\]
где \(p>2\) и \(1<q<2\). В дискретном времени возможны
перескакивание через решение и колебания. Поэтому непрерывная
фиксированно-временная устойчивость не означает точного достижения за
фиксированное число шагов. Для дискретной схемы требуется отдельный выбор достаточно малого шага; универсальной границы на \(\eta\) для общего нелинейного случая нет.

\section{Фиксированно-временной ньютоновский поток}

Пусть
\[
    f\in C^2(\RR^n),
    \qquad
    \nabla^2f(x)\succ0
\]
во всех рассматриваемых точках, а минимум \(x^\star\) существует.
Тогда \(f\) строго выпукла, минимум единственен, а Гессиан обратим.
Обозначим
\[
    g(x)=\nabla f(x),
    \qquad
    H(x)=\nabla^2f(x).
\]
Обычный ньютоновский поток имеет вид
\[
    \dot x=-H(x)^{-1}g(x).
\]
Для
\[
    V(x)=\frac12\lVert g(x)\rVert^2
\]
получаем
\[
    \dot V
    =
    g(x)^{\mathsf T}H(x)\dot x
    =
    -\lVert g(x)\rVert^2
    =
    -2V.
\]
Следовательно,
\[
    \lVert\nabla f(x(t))\rVert
    =
    e^{-(t-t_0)}
    \lVert\nabla f(x(t_0))\rVert.
\]

Выберем параметры
\[
    p>2,
    \qquad
    1<q<2,
\]
и положим
\[
    \rho_p=\frac{p-2}{p-1},
    \qquad
    \rho_q=\frac{q-2}{q-1}.
\]
Масштабированный поток задаётся уравнением
\[
\begin{aligned}
    \dot x
    ={}&
    -H(x)^{-1}
    \frac{g(x)}{\lVert g(x)\rVert^{\rho_p}}
    \\
    &-
    H(x)^{-1}
    \frac{g(x)}{\lVert g(x)\rVert^{\rho_q}},
\end{aligned}
\]
а при \(g(x)=0\) полагается \(\dot x=0\). Обычный и
масштабированный ньютоновские потоки имеют одинаковые геометрические
траектории. Они в общем случае не совпадают с траекториями обычного
градиентного потока.

\section{Функция Ляпунова и оценка времени}

Поскольку строгая выпуклость и существование минимума дают
\[
    g(x)=0
    \quad\Longleftrightarrow\quad
    x=x^\star,
\]
функция
\[
    V(x)=\frac12\lVert g(x)\rVert^2
\]
положительно определена относительно \(x^\star\). Вдоль
масштабированного ньютоновского потока
\[
\begin{aligned}
    \dot V
    ={}&
    -\lVert g(x)\rVert^{2-\rho_p}
    \\
    &-
    \lVert g(x)\rVert^{2-\rho_q}.
\end{aligned}
\]
Положим
\[
    \alpha_1=\frac{p}{2(p-1)},
    \qquad
    \alpha_2=\frac{q}{2(q-1)}.
\]
Так как \(\lVert g(x)\rVert^2=2V(x)\), имеем
\[
    \dot V
    =
    -2^{\alpha_1}V^{\alpha_1}
    -
    2^{\alpha_2}V^{\alpha_2}.
\]
При \(p>2\) выполняется \(0<\alpha_1<1\), а при
\(1<q<2\) --- \(\alpha_2>1\).

\begin{theorem}[Фиксированное время ньютоновского потока]
Пусть \(f\in C^2(\RR^n)\), минимум \(x^\star\) существует,
\(\nabla^2f(x)\succ0\) вдоль рассматриваемых траекторий,
\(p>2\) и \(1<q<2\). Если соответствующее решение существует на
всём промежутке до достижения равновесия, то
\(x^\star\) фиксированно-временно устойчиво и
\[
    T(x_0)
    \le
    \frac1{2^{\alpha_1}(1-\alpha_1)}
    +
    \frac1{2^{\alpha_2}(\alpha_2-1)}.
\]
\end{theorem}

В оценке не появляется константа сильной выпуклости: в производной
\(V\) Гессиан сокращается с обратным Гессианом. Однако условие
\(H(x)\succ0\) не даёт равномерной оценки
\[
    H(x)\succeq\mu I.
\]
Поэтому \(\lVert H(x)^{-1}\rVert\) может быть большой в областях
малой кривизны. Полный результат требует контролировать существование
траектории, отсутствие ухода на бесконечность, корректность поля при
малых собственных значениях Гессиана и продолжение решения после
достижения равновесия.

Можно ввести коэффициенты \(a,b>0\):
\[
\begin{aligned}
    \dot x
    ={}&
    -aH(x)^{-1}
    \frac{g(x)}{\lVert g(x)\rVert^{\rho_p}}
    \\
    &-
    bH(x)^{-1}
    \frac{g(x)}{\lVert g(x)\rVert^{\rho_q}}.
\end{aligned}
\]
Тогда
\[
    T_{\max}
    \le
    \frac1{a2^{\alpha_1}(1-\alpha_1)}
    +
    \frac1{b2^{\alpha_2}(\alpha_2-1)}.
\]
Коэффициенты можно выбрать так, чтобы эта верхняя граница не
превышала заранее заданного времени. Параметры \(p,q\) задают степени,
а \(a,b\) непосредственно масштабируют длительность движения.

\section{PL-неравенство и квадратичный рост}

Пусть дифференцируемая функция достигает значения \(f^\star\) и
удовлетворяет неравенству Поляка--Лоясиевича
\[
    \frac12\lVert\nabla f(x)\rVert^2
    \ge
    \mu\bigl(f(x)-f^\star\bigr),
    \qquad
    \mu>0.
\]
Выпуклость здесь не предполагается. Из условия следует
\[
    \nabla f(x)=0
    \quad\Longrightarrow\quad
    f(x)=f^\star,
\]
то есть неоптимальных стационарных точек нет. Множество минимумов
\[
    \mathcal X^\star
    =
    \operatorname*{argmin}_{x\in\RR^n}f(x)
\]
при этом может содержать несколько точек.

\begin{definition}[Квадратичный рост]
Функция удовлетворяет условию квадратичного роста с константой
\(\mu>0\), если
\[
    f(x)-f^\star
    \ge
    \frac\mu2
    \operatorname{dist}(x,\mathcal X^\star)^2,
\]
где
\[
    \operatorname{dist}(x,\mathcal X^\star)
    =
    \inf_{z\in\mathcal X^\star}\lVert x-z\rVert.
\]
\end{definition}

\begin{theorem}[PL-неравенство влечёт квадратичный рост]
Пусть \(f\) дифференцируема, достигает минимума и удовлетворяет
PL-неравенству с константой \(\mu>0\). При условиях регулярности,
обеспечивающих существование используемого ниже градиентного потока до
достижения \(\mathcal X^\star\), выполняется
\[
    f(x)-f^\star
    \ge
    \frac\mu2
    \operatorname{dist}(x,\mathcal X^\star)^2
\]
для всех \(x\).
\end{theorem}

Таким образом, малая ошибка по значению возможна только вблизи
множества глобальных решений. При единственном минимизаторе условие
принимает привычный вид
\[
    f(x)-f^\star
    \ge
    \frac\mu2\lVert x-x^\star\rVert^2,
\]
но ни PL-условие, ни квадратичный рост сами по себе не требуют
выпуклости.

\section{Доказательство квадратичного роста}

Возьмём \(x_0\notin\mathcal X^\star\) и определим
\[
    G(x)=\sqrt{f(x)-f^\star}.
\]
Вне множества решений
\[
    \nabla G(x)
    =
    \frac{\nabla f(x)}
    {2\sqrt{f(x)-f^\star}}.
\]
Поэтому
\[
    \lVert\nabla G(x)\rVert^2
    =
    \frac{\lVert\nabla f(x)\rVert^2}
    {4(f(x)-f^\star)}
    \ge
    \frac\mu2.
\]
Рассмотрим поток
\[
    \dot x=-\nabla G(x)
\]
до первого момента достижения \(\mathcal X^\star\). Вдоль него
\[
\begin{aligned}
    \frac{\mathrm d}{\mathrm dt}G(x(t))
    &=
    -\lVert\nabla G(x(t))\rVert^2
    \\
    &\le
    -\frac\mu2.
\end{aligned}
\]
Следовательно, множество решений достигается не позднее времени
\[
    T\le\frac{2G(x_0)}\mu.
\]
Обозначим конечную точку через \(x_T\in\mathcal X^\star\).

С другой стороны,
\[
\begin{aligned}
    G(x_0)
    &=
    \int_0^T
    \lVert\nabla G(x(t))\rVert^2\,\mathrm dt
    \\
    &\ge
    \sqrt{\frac\mu2}
    \int_0^T
    \lVert\nabla G(x(t))\rVert\,\mathrm dt.
\end{aligned}
\]
Поскольку
\[
    \lVert\dot x(t)\rVert
    =
    \lVert\nabla G(x(t))\rVert,
\]
последний интеграл является длиной траектории и не меньше расстояния
между её концами:
\[
    \int_0^T
    \lVert\dot x(t)\rVert\,\mathrm dt
    \ge
    \lVert x_0-x_T\rVert
    \ge
    \operatorname{dist}(x_0,\mathcal X^\star).
\]
Значит,
\[
    G(x_0)
    \ge
    \sqrt{\frac\mu2}
    \operatorname{dist}(x_0,\mathcal X^\star).
\]
После возведения в квадрат получаем
\[
    f(x_0)-f^\star
    \ge
    \frac\mu2
    \operatorname{dist}(x_0,\mathcal X^\star)^2.
\]
Поскольку начальная точка произвольна, теорема доказана.

\section{Следствия, примеры и границы}

Из квадратичного роста следует
\[
    f(x)-f^\star\le\varepsilon
    \quad\Longrightarrow\quad
    \operatorname{dist}(x,\mathcal X^\star)
    \le
    \sqrt{\frac{2\varepsilon}{\mu}}.
\]
Из PL-неравенства получаем
\[
    f(x)-f^\star
    \le
    \frac1{2\mu}
    \lVert\nabla f(x)\rVert^2.
\]
Объединение двух оценок даёт
\[
    \operatorname{dist}(x,\mathcal X^\star)
    \le
    \frac1\mu
    \lVert\nabla f(x)\rVert.
\]
Поэтому малая норма градиента у PL-функции означает близость к
множеству глобальных решений, а не только почти стационарность.

PL-условие не гарантирует единственности минимума. Для
\[
    f(x_1,x_2)=\frac12x_1^2
\]
имеем
\[
    \frac12\lVert\nabla f(x)\rVert^2
    =
    f(x)-f^\star,
\]
но
\[
    \mathcal X^\star
    =
    \{(0,s):s\in\RR\}.
\]
При этом
\[
    f(x)-f^\star
    =
    \frac12
    \operatorname{dist}(x,\mathcal X^\star)^2.
\]
Следовательно, утверждение, будто PL-функция либо постоянна, либо
имеет единственный минимум, неверно.

PL-свойство исключает неоптимальные стационарные точки; инвексность является отдельным структурным условием. Функция
\[
    f(x)=x^2+3\sin^2x
\]
иллюстрирует невыпуклую геометрию; без отдельной проверки PL-неравенства она не используется как доказательный пример.

Для дифференцируемых функций сохраняется цепочка
\[
    \text{сильная выпуклость}
    \Longrightarrow
    \text{PL}
    \Longrightarrow
    \text{квадратичный рост},
\]
при указанных предпосылках существования минимума и регулярности.
Обратные импликации в общем случае неверны.

Ньютоновская конструкция требует решения системы
\[
    H(x)d=-g(x),
\]
а не явного вычисления обратной матрицы, но такая операция всё равно
может быть дорогой. PL-результат требует
\(\mathcal X^\star\ne\varnothing\), а доказательство через \(G\) ---
существования траектории до достижения множества решений.
Непрерывные фиксированно-временные оценки не переносятся автоматически
на методы с конечным шагом.
